\section{Measurement Methods}
\label{sec:methods}

This section defines what each stage measures. Throughout, $x\in\mathbb{R}^{T}$ is a context of length
$T$ (512 in the case study), $y\in\mathbb{R}^{H}$ the true continuation ($H=64$), $\hat y$ a model's point
forecast and $h^{(\ell)}$ the residual stream after block $\ell$.

\subsection{Behavior and cost (L0)}
\label{sec:l0}

Accuracy is scored per series with MASE \citep{hyndman2006mase},
\begin{equation}
  \mathrm{MASE}(\hat y, y) \;=\; \frac{\tfrac1H\sum_{t=1}^{H}|y_t-\hat y_t|}
       {\max\!\big(\tfrac{1}{T-m}\sum_{t=m+1}^{T}|x_t-x_{t-m}|,\;\epsilon_x\big)},
\end{equation}
with the seasonal lag $m$ taken from ground truth where available, and a scale floor $\epsilon_x$
because on intermittent series the naive denominator collapses and MASE spuriously improves
(\fid{MN-23}). sMAPE, pinball loss \citep{gneiting2007scoring} and interval calibration are reported
alongside. Model pairs are compared with a paired bootstrap over series, per family and overall, with
Holm correction across all pairs and families \citep{holm1979}. Each model's repeat-run \emph{noise floor}
(two identical runs with different seeds) is measured so that a difference can be read in floor units;
a deterministic model has a floor of zero, which the report states rather than using as a referee. The
\texttt{budget} stage measures parameters by role, FLOPs with PyTorch's \texttt{FlopCounterMode},
latency, memory and the share of forecast FLOPs spent inside captured blocks (Table~\ref{tab:models}).
The \texttt{frontend} stage tests the tokenizer before any layer runs: scale equivariance (forecasting
$c\,x$ should give $c\,\hat y$), NaN handling and quantization resolution.

\subsection{Geometry and linear translatability (L1, L2)}
\label{sec:l1l2}

For aligned activations $A\in\mathbb{R}^{n\times d_a}$ and $B\in\mathbb{R}^{n\times d_b}$ over the same
$n$ series (mean-pooled over windows, then centered), linear CKA \citep{kornblith2019cka} is
\begin{equation}
  \mathrm{CKA}(A,B) \;=\; \frac{\|B^{\top}A\|_F^2}{\|A^{\top}A\|_F\,\|B^{\top}B\|_F}.
\end{equation}
It is computed for every layer pair of every model pair, with a cluster bootstrap over series and a
shuffled-series null that breaks correspondence while preserving each model's marginal geometry. L2 fits
ridge maps between layers and reports only the \emph{gain} of the stitched prediction over an
input-feature baseline, never a raw $R^2$, because a large part of any layer is linearly predictable from
the input alone. Activation clusterings are compared across models with adjusted mutual information
\citep{vinh2010ami}. These are \ev{geometric} and \ev{linear-translatable} claims; they say nothing about
what the model uses.

\subsection{Perturbation and patching (L3)}
\label{sec:l3}

For each corruption $c$ of the nine in Section~\ref{sec:benchmark}, L3 records a sensitivity
\emph{fingerprint}, the relative activation change $\|h^{(\ell)}(c(x))-h^{(\ell)}(x)\|/\|h^{(\ell)}(x)\|$ per
layer, and a within-model \emph{restoration curve}: the fraction of the forecast damage undone by
patching the clean activations of layer $\ell$ (whole layer, or one window at a time) into the corrupted
forward pass. Two models are compared by rank-correlating their fingerprints on the relative depth range
they share, never by transplanting activations between models. Restoration is \ev{causal within-model};
agreement between models is correlational.

\subsection{Depth: lens, internals and coverage}
\label{sec:depth}

A skip lens decodes the forecast from each layer through the model's own head, and a tuned ridge lens
\citep{belrose2023tuned} fits a per-layer readout; the crystallization depth is the shallowest layer
whose lens forecast is within tolerance of the final one. \texttt{internals} reports per-layer effective
dimension (participation ratio), decodability of benchmark families and CKA to the input.
\texttt{layer\_screen} selects the layers that deserve expensive analysis with a cheap proxy
(\texttt{work\_bend}). Two design points qualify every depth claim. First, the depth axis is the fraction of
the \emph{whole} stack's blocks, including uncaptured parts, so Chronos-T5's last captured encoder block
sits at 0.478 rather than 1.0 (\fid{DE-04}). Second, the lens is withheld for a model whose head reads
positions the patch never writes (\texttt{forecast\_reads\_patched\_positions}), because such a lens returns
a perfectly flat and meaningless curve (\fid{DE-03}).

\subsection{Causal concepts}
\label{sec:concepts}

\paragraph{Sparse dictionaries.} For each target layer, \tool{} trains a TopK sparse autoencoder
\citep{makhzani2013ksparse,gao2025topk} on token-level activations,
\begin{equation}
  z(h) = \operatorname{TopK}\!\big(W_e(h-b_d)+b_e\big),\qquad \hat h = W_d\,z(h) + b_d ,
\end{equation}
with the dictionary size chosen by a search, an auxiliary loss that revives dead atoms and a minimum
number of training steps (\fid{MP-03}). Fidelity, dead rate and forecast preservation are reported per
target. Because a single SAE seed is not reproducible \citep{paulo2025seeds}, $n_{\text{seeds}}$ replicate
dictionaries are trained at the primary's size to measure a within-model reproducibility ceiling
(\fid{MP-04}).

\paragraph{The ablation battery.} A feature $i$ is tested on the $k$ series on which it fires most. For
those series, the layer's activations are replaced by the SAE reconstruction with $z_i$ set to zero, and the
forecast is compared with the forecast under the full reconstruction:
\begin{equation}
  \Delta_i^{(c)}(x) \;=\; g_c\!\big(f(x;\,h^{(\ell)}\!\leftarrow\!\hat h_{-i})\big) - g_c\!\big(f(x;\,h^{(\ell)}\!\leftarrow\!\hat h)\big),
  \label{eq:ablation}
\end{equation}
where $f(x;\,h\!\leftarrow\!\cdot)$ is the forecast with layer $\ell$ patched and $g_c$ is one of nine
forecast \emph{channels}: trend, seasonal strength, spectral centroid, level, dispersion, near- and
far-horizon shape distance, MASE and flatness. Using the full reconstruction, not the clean forecast, as the
baseline keeps the dictionary's own reconstruction error from being charged to the feature. The
\emph{null} removes a random unit direction in code space, scaled to the mean magnitude the ablations
remove, on the \emph{same} series; the feature clears channel $c$ if its mean $|\Delta_i^{(c)}|$ exceeds
the 95th percentile of the null draws restricted to its own rows. The row matching matters: a feature
whose top series happen to sit in a high-response regime would otherwise clear a pooled null on that
alone. The channel vector used downstream is $v_i^{(c)} = \overline{\Delta_i^{(c)}}/q_{95}^{(c)}$, the
signed mean effect in null units. The near- and far-horizon shape channels are distances ($\ge 0$) and are
treated as undirected everywhere, including in generated text. Before any feature is scored, a
\emph{reach probe} requires that patching the layer's own clean activations into itself changes the
forecast by exactly $0.0$ and that patching another layer's activations changes it by a relative amount
above $10^{-3}$ (Figure~\ref{fig:battery}). A target that fails the probe is withheld with its reason,
because a table of zeros from a patch that never lands reads as ``these features do not matter''.

\begin{figure}[t]
\centering
\begin{tikzpicture}[
  box/.style={draw, rounded corners=2pt, align=center, font=\scriptsize, minimum height=9mm, fill=white, text width=21mm},
  gate/.style={draw, diamond, aspect=2.0, align=center, font=\scriptsize, inner sep=0.5pt, fill=soft, text width=15mm},
  arr/.style={-{Stealth[length=1.6mm]}, gray!70!black},
  lab/.style={font=\scriptsize\itshape, text=gray!70!black}]
  \node[box] (tgt) at (0,0) {target layer $\ell$\\+ trained SAE};
  \node[gate] (self) at (3.5,0) {self-patch\\$=0.0$?};
  \node[gate] (cross) at (7.0,0) {cross-patch rel.\\$>10^{-3}$?};
  \node[box] (rows) at (10.6,0) {feature $i$: its\\top-$k$ firing series};
  \node[box, fill=bad!8] (wh) at (5.25,1.7) {target withheld,\\reason rendered};
  \node[box] (abl) at (10.6,-2.2) {forecast with $\hat h_{-i}$\\vs.\ with $\hat h$\\(same seed)};
  \node[box] (null) at (7.0,-2.2) {row-matched null:\\random code direction,\\matched magnitude};
  \node[box] (chan) at (3.5,-2.2) {each of 9 channels:\\$|\overline{\Delta}|>q_{95}$ of null?};
  \node[box, fill=accent!12] (vec) at (0,-2.2) {causal feature\\$v_i=\overline{\Delta}/q_{95}$};
  \draw[arr] (tgt) -- (self);
  \draw[arr] (self) -- node[above,lab]{yes} (cross);
  \draw[arr] (cross) -- node[above,lab]{yes} (rows);
  \draw[arr] (self.north) |- node[left,lab,pos=0.3]{no} (wh.west);
  \draw[arr] (cross.north) |- node[right,lab,pos=0.3]{no} (wh.east);
  \draw[arr] (rows) -- (abl);
  \draw[arr] (abl) -- (null);
  \draw[arr] (null) -- (chan);
  \draw[arr] (chan) -- node[below,lab]{$\ge$1 clears} (vec);
\end{tikzpicture}
\caption{The ablation battery. A target is scored only after both reach controls pass; each feature is
ablated on its own top-firing series and compared, channel by channel, with a random-direction null on
exactly those series. Only features that clear at least one channel enter the concept analyses.}
\label{fig:battery}
\end{figure}

\paragraph{From features to concepts.} Causal features are grouped at three granularities, each
answering a different question. \emph{Per-target concepts} cluster one layer's feature vectors with
$k$-means, choosing $k$ by silhouette \citep{rousseeuw1987silhouette} among values whose smallest cluster
has at least three members; a target with no admissible $k$ is reported as \emph{non-modular}. (Without
the minimum-size constraint, silhouette rewards splitting off a single outlier; \fid{SH-20}.) The
\emph{concept atlas} pools causal features from every model and layer and cuts a complete-linkage
dendrogram at cosine 0.9, which guarantees that every pair of members has cosine at least 0.9; an atlas
concept is a shared \emph{causal-effect profile}, not a shared input. \emph{Families} are a coarser,
readable tiling: average linkage plus nearest-centroid assignment at cosine 0.5, with the number of
families chosen by a parsimony rule and tested against a column-shuffle null. Families are titled
deterministically from their directed channels (Figure~\ref{fig:families}).

\subsection{Two notions of sharing}
\label{sec:sharing}

Whether a concept is ``shared'' between models is two separate questions, and conflating them produces
wrong answers in both directions (Section~\ref{sec:controls}). \tool{} answers them separately and then
combines them in an evidence ladder (Figure~\ref{fig:ladder}).

\paragraph{Same inputs.} \emph{Concept transfer} asks whether some feature in another model separates a
concept's top-$k$ firing series from the rest, scored by AUC, against a null of random series subsets
matched on benchmark strata (family and archetype); the test is run in both directions and a transfer is
\emph{reciprocal} if both legs clear. At atlas scale, $p$-values are corrected by Benjamini--Hochberg
\citep{benjamini1995fdr} per ordered model pair and leg. When concept profiles compare what two concept
parts fire on, input agreement is measured as the overlap of their top-$k$ series (hypergeometric test
with BH), never as a whole-series correlation, because two sparse features can have $\rho=0.80$ with no
shared top series (\fid{MN-19}).

\paragraph{Same effect on the same inputs.} \emph{Shared-input causal agreement} takes every
reciprocal-FDR transfer and ablates the feature (or concept part) on each side on the \emph{same} series
set $U$. Each side must first clear its own random-direction null on $U$ to be scorable. Two statistics,
level concordance and the cosine of the shape effects, must then beat both sides'
\emph{activation-matched random-feature floors}: sets of equally active features whose ablation defines
what agreement looks like by chance. A pair is ``same causal effect'' if both statistics beat both floors,
``level only'' or ``shape only'' if one does, and ``acts differently'' only if a statistic falls below the
floors' 5th percentile; failing to beat the 95th percentile is ``no specific agreement''. The distinction
matters: absence of agreement is not evidence of disagreement (\fid{MN-18}).

\paragraph{Sharing classes.} Each atlas concept then receives a \texttt{sharing\_class}: \emph{shared}
(same effect and same inputs across models), \emph{partially shared}, \emph{convergent} (same effect,
different inputs) or \emph{single-model}. A concept part is flagged provenance-driven when at least 80\% of
its top series come from one real-derived generator.

\begin{figure}[t]
\centering
\begin{tikzpicture}[
  rung/.style={draw, rounded corners=2pt, font=\scriptsize, align=left, text width=0.78\textwidth, minimum height=6.5mm, fill=white},
  tag/.style={font=\scriptsize\bfseries, anchor=east},
  arr/.style={-{Stealth[length=1.5mm]}, gray!60!black}]
  \node[rung] (r1) {\textbf{Same forecast effect:} atlas groups causal features of $\ge$2 models at pairwise cosine $\ge$0.9 \hfill \textit{causal, compared}};
  \node[rung, above=1.2mm of r1] (r2) {\textbf{Same inputs:} the parts' top-$k$ series overlap beyond chance (hypergeometric, BH) \hfill \textit{correlational}};
  \node[rung, above=1.2mm of r2] (r3) {\textbf{Reproducible:} concept survives independent SAE seeds (within-model ceiling) \hfill \textit{stability}};
  \node[rung, above=1.2mm of r3] (r4) {\textbf{Other dictionaries select the same inputs:} reciprocal transfer vs.\ stratum-matched null \hfill \textit{correlational}};
  \node[rung, above=1.2mm of r4] (r5) {\textbf{Same causal effect on the same inputs:} joint ablation on $U$ beats matched floors \hfill \textit{causal on shared inputs}};
  \node[rung, above=1.2mm of r5, fill=accent!10] (r6) {\textbf{Confirmed on private data:} registered claim confirmed once on the sealed split \hfill \textit{held-out confirmed}};
  \foreach \i in {1,...,6}{ \node[tag] at ($(r\i.west)+(-1mm,0)$) {L\i}; }
\end{tikzpicture}
\caption{The concept evidence ladder used by the report (computed by \texttt{concept\_verdicts} in \texttt{report/derived.py}).
Every causal feature has already passed the ablation battery (Figure~\ref{fig:battery}). Each rung is a
separate test with its own null, and the report states the highest rung each concept reaches together with
the rule that produced its verdict. ``Shared'' in this paper means the \texttt{shared} sharing class: same
effect (L1) and same inputs (L2).}
\label{fig:ladder}
\end{figure}
